\clearpage
\SourceLeaf{169}%
\pdfbookmark[0]{Göttingische gelehrte Anzeigen. 1817 März 10.}{notice1817march10}
\begin{center}
\rule{.82\linewidth}{.35pt}\par
Göttingische gelehrte Anzeigen.\quad 1817 März 10.\par
\rule{.82\linewidth}{.35pt}
\end{center}
Am 10. Februar wurde der Königl. Societät von Hrn. Hofr. \textsc{Gauss} eine Vorlesung eingereicht, überschrieben:
\begin{center}\itshape
Theorematis fundamentalis in doctrina de residuis quadraticis demonstrationes et ampliationes novae.
\end{center}

Es ist eine Eigenthümlichkeit der höhern Arithmetik, dass so viele ihrer schönsten Lehrsätze mit grösster Leichtigkeit durch Induction entdeckt werden können, deren Beweise jedoch nichts weniger als nahe liegen, sondern oft erst nach vielen vergeblichen Versuchen mit Hülfe tiefeindringender Untersuchungen und glücklicher Combinationen gefunden werden. Dies merkwürdige Phänomen entspringt aus der oft wunderbaren Verkettung der verschiedenartigen Lehren in jenem Theile der Mathematik, und eben daher kommt es, dass häufig solche Lehrsätze, von denen anfangs ein Beweis Jahre lang vergeblich gesucht war, späterhin sich auf mehreren ganz verschiedenen Wegen beweisen lassen. Sobald ein neuer Lehrsatz durch Induction entdeckt ist, hat man die Auffindung \emph{irgend eines} Beweises freilich als das erste Erforderniss zu betrachten: allein nachdem ein solcher geglückt ist, darf man in der höhern Arithmetik die Untersuchung nicht immer als abgeschlossen und die Aufspürung anderer Beweise als überflüssigen Luxus ansehen. Denn theils kommt man gewöhnlich auf die schönsten und einfachsten
\SourceLeaf{170}%
Beweise nicht zuerst, und dann ist gerade die Einsicht in die wunderbare Verkettung der Wahrheit der höhern Arithmetik dasjenige, was einen Hauptreiz dieses Studiums ausmacht, und nicht selten wiederum zur Entdeckung neuer Wahrheiten führt. Aus diesen Gründen ist hier die Auffindung neuer Beweise für schon bekannte Wahrheiten öfters für wenigstens eben so wichtig anzusehen, als die Entdeckung der Wahrheit selbst. Kennern der höhern Arithmetik sind diese Betrachtungen nicht neu; man weiss, dass ein grosser Theil von \textsc{Eulers} Verdiensten um dieselbe in der Auffindung von Beweisen für Lehrsätze besteht, die schon von \textsc{Fermat} wie es scheint durch Induction gefunden waren.

Die Lehre von den quadratischen Resten gibt einen einleuchtenden Beleg zu dem vorhin Gesagten. Sie beruhet hauptsächlich auf dem sogenannten Fundamental-Theorem, welches darin besteht, dass die wechselseitigen Relationen zweier (ungeraden positiven) Primzahlen zu einander (in sofern der eine quadratischer Rest oder Nichtrest der andern ist) einerlei sind, so oft eine der Primzahlen oder beide unter der Form $4k+1$ stehen, entgegengesetzt aber, so oft beide Primzahlen von der Form $4k+3$ sind. Für solche Leser, die mit der höhern Arithmetik weniger bekannt sind, erinnern wir, dass eine ganze Zahl quadratischer Rest einer andern heisst, wenn die erstere um ein Vielfaches der andern vermehrt ein Quadrat geben kann; Nichtrest hingegen, wenn dies nicht möglich ist. Die Geschichte dieses schönen durch Induction äusserst leicht zu findenden Lehrsatzes wollen wir hier nicht vollständig wiederholen, sondern nur bemerken, dass der Verfasser vorliegender Abhandlung, nach Anfangs ziemlich lange vergeblich angestellten Untersuchungen, nach und nach bereits vier unter sich ganz verschiedene Beweise gegeben hat, wovon zwei in den \emph{Disquisitionibus Arithmeticis} enthalten sind, der dritte den Gegenstand einer eigenen Abhandlung im sechzehnten Bande der Commentationen ausmacht, und der vierte in eine Abhandlung \emph{summatio quarumdam serierum singularium} im ersten Bande der \emph{Commentationes recentiores} verwebt ist; über diese beiden Abhandlungen kann man unsere Anzeigen 1808. Mai 12 und Sept. 19 nachsehen, wo auch vollständigere geschichtliche Nachweisungen befindlich sind. Dass der Verf. bei diesen vier Beweisen, ungeachtet jeder derselben für sich in Rücksicht auf Strenge nichts zu wünschen übrig lässt, noch nicht stehen geblieben ist, bedarf zwar bei den Freunden der höhern Arithmetik keiner Rechtfertigung; indessen würde er doch wahrscheinlich sich nicht so eifrig bemüht haben, jenen Beweisen noch andere hinzuzufügen, wenn
\SourceLeaf{171}%
nicht ein besonderer Umstand ihn dazu veranlasst hätte, der hier erwähnt werden muss. Seit dem Jahre 1805 hatte er nemlich angefangen, sich mit den Theorien der cubischen und biquadratischen Reste zu beschäftigen, welche noch weit reichhaltiger und interessanter sind, als die Theorie der quadratischen Reste. Es zeigten sich bei jenen Untersuchungen dieselben Erscheinungen wie bei der letztern, nur gleichsam mit vergrössertem Massstabe. Durch Induction, sobald nur der rechte Weg dazu eingeschlagen war, fanden sich sogleich eine Anzahl höchst einfacher Theoreme, die jene Theorien ganz erschöpfen, mit den für die quadratischen Reste geltenden Lehrsätzen eine überraschende Aehnlichkeit haben, und namentlich auch zu dem Fundamentaltheorem das Gegenstück darbieten. Allein die Schwierigkeiten, für jene Lehrsätze ganz befriedigende Beweise zu finden, zeigten sich hier noch viel grösser, und erst nach vielen, eine ziemliche Reihe von Jahren hindurch fortgesetzten Versuchen ist es dem Verfasser endlich gelungen, sein Ziel zu erreichen. Die grosse Analogie der Lehrsätze selbst, bei den quadratischen und bei den höhern Resten, liess vermuthen, dass es auch analoge Beweise für jene und diese geben müsse; allein die zuerst für die quadratischen Reste gefundenen Beweisarten vertrugen gar keine Anwendung auf die höhern Reste, und gerade dieser Umstand war der Beweggrund, für jene immer noch andere neue Beweise aufzusuchen. Der Verf. wünscht daher, dass man die vorliegende Abhandlung, die für die Theorie der quadratischen Reste noch einige neue Hülfsquellen eröffnet, als Vorläuferin der Theorie der cubischen und biquadratischen Reste betrachte, die er in Zukunft bekannt zu machen denkt, und die zu den schwierigsten Gegenständen der höhern Arithmetik gehören.

Die gegenwärtige Abhandlung besteht aus dreien von einander unabhängigen Theilen. Sie enthält nemlich den fünften und sechsten Beweis des Fundamental-Theorems und eine neue, mit dem dritten Beweise zusammenhängende Methode,\textperiodcentered\ zu entscheiden, ob eine vorgegebene ganze Zahl von einer gegebenen Primzahl quadratischer Rest oder Nichtrest sei. Unter den vier ersten Beweisen war der dritte unstreitig derjenige, der die grösste Einfachheit mit Unabhängigkeit von fremdartigen Untersuchungen vereinigte, daher ihn auch \textsc{Legendre} in die neue Ausgabe seines \emph{Essai d'une théorie des nombres} aufgenommen hat. Der \emph{fünfte} Beweis scheint dem dritten in beiden Hinsichten wenigstens gleich zu kommen. Beide Beweise haben insofern einige Verwandtschaft, dass sie von einem und demselben Lehrsatze ausgehen, sind aber bei der weitern Ausführung völlig von \SourceLeaf{172}%
einander verschieden. Dieser Lehrsatz besteht in Folgendem: Wenn $m$ eine (positive ungerade) Primzahl; $M$ eine ganze durch $m$ nicht theilbare Zahl bedeutet, wenn ferner unter den Resten, die aus der Division der Producte
\[
M,\ 2M,\ 3M,\ 4M\mathinner{.\,.\,.\,.\,.\,.}\tfrac12(m-1)M
\]
durch $m$ entstehen, die Anzahl derjenigen, die grösser als $\frac12m$ sind, durch $n$ bezeichnet wird, so ist $M$ quadratischer Rest oder Nichtrest von $m$, jenachdem $n$ gerade oder ungerade ist. Um nun zu dem Beweise des Fundamentallehrsatzes zu gelangen, wird angenommen, dass auch $M$ eine ungerade positive Primzahl und $N$ in Beziehung auf $M$ und $m$ dasselbe bedeutet, was $n$ in Beziehung auf $m$ und $M$ ausdrückt, so dass $N$ gerade oder ungerade entscheidet, ob $m$ quadratischer Rest oder Nichtrest von $M$ ist. Durch eine sehr kurze Reihe von Schlüssen zeigt der Verfasser, dass die Anzahl aller positiven ganzen Zahlen, die zugleich kleiner als $\frac12mM$ sind, mit $m$ dividirt einen Rest kleiner als $\frac12m$, und mit $M$ dividirt einen Rest kleiner als $\frac12M$ geben,
\[
=\tfrac18(m-1)(M-1)+\tfrac12n+\tfrac12N
\]
und folglich allemal
\[
\tfrac14(m-1)(M-1)+n+N
\]
eine gerade Zahl sei. So oft also wenigstens eine der Zahlen $m,M$ von der Form $4k+1$ ist, mithin $\frac14(m-1)(M-1)$ gerade, wird auch $n+N$ gerade sein, folglich entweder $n$ und $N$ beide gerade, oder beide ungerade. Wenn hingegen sowohl $m$ als $M$ von der Form $4k+3$ ist, wird nothwendig $n+N$ ungerade, folglich eine der Zahlen $n,N$ gerade, die andere ungerade sein. Hieraus folgt in Verbindung mit obigem Lehrsatze das Fundamental-Theorem von selbst.

Der \emph{sechste} Beweis ist zwar von gleicher Kürze und Concinnität wie der fünfte, beruhet aber doch auf etwas künstlichern Combinationen. Der beschränkte Raum dieser Blätter erlaubt nur, mit Uebergehung des Einzelnen, hier das Hauptmoment zu berühren. Es bezeichnen
\begin{description}
\item[$p,q$] zwei (ungleiche positive ungerade) Primzahlen,
\item[$\alpha$] eine sogenannte \emph{radix primitiva} für den Modulus $p$, d. i. eine durch $p$ nicht theilbare (hier positive) ganze Zahl von der Art, dass keine niedrigere Potenz als $\alpha^{p-1}$ nach dem Modulus $p$ der Einheit congruent wird
\item[$x$] eine unbestimmte Grösse
\SourceLeaf{173}%
\item[$\xi$] die Function
\[
x-x^\alpha+x^\zeta-x^\eta+x^\theta-\text{etc.}-x^\lambda
\]
wo (des bequemern Drucks wegen) $\zeta,\eta,\theta\ldots\lambda$ statt der Zahlen $\alpha\alpha,\allowbreak \alpha^3,\allowbreak \alpha^4\ldots\allowbreak \alpha^{p-2}$ gesetzt sind;
\item[$\varepsilon$] die Einheit, positiv genommen, wenn $p$ von der Form $4k+1$, negativ, wenn $p$ von der Form $4k+3$ ist;
\item[$\delta$] die Einheit, positiv genommen, wenn wenigstens eine der Zahlen $p,q$ von der Form $4k+1$ ist, negativ, wenn beide von der Form $4k+3$ sind;
\item[$\gamma$] die Einheit, positiv genommen, wenn $q$ ein quadratischer Rest von $p$ ist, negativ, wenn $q$ quadratischer Nichtrest von $p$ ist;
\item[$\beta$] die Einheit, positiv genommen, wenn $p$ ein quadratischer Rest von $q$, negativ, wenn $p$ quadratischer Nichtrest von $q$ ist.
\end{description}

Nach diesen Vorbereitungen folgt leicht aus dem 51. Art. der \emph{Disquisitiones Arithmeticae}, dass die Function
\[
\xi^q-x^q+x^{q\alpha}-x^{q\zeta}+x^{q\eta}-x^{q\theta}+\text{etc.}+x^{q\lambda}
\]
entwickelt lauter durch $q$ theilbare Coëfficienten bekommt, und daher, wenn diese Function $=qX$ gesetzt wird, $X$ eine auch in Beziehung auf die Coëfficienten \emph{ganze} Function werde. Durch Schlüsse, in die näher einzugehen hier zu weitläufig sein würde, wird in der Abhandlung bewiesen, dass die Function $qX\xi$ mit $x^{p-1}+x^{p-2}+x^{p-3}+x^{p-4}+\text{etc.}+x+1$ dividirt, den Rest
\[
\varepsilon p(\delta p^{\frac12(q-1)}-\gamma)
\]
gibt, daher aus der Division der Function $X\xi$ mit demselben Divisor der Rest
\[
\frac{\varepsilon p(\delta p^{\frac12(q-1)}-\gamma)}{q}
\]
hervorgehen wird. Diese Grösse muss daher nothwendig eine ganze Zahl sein, woraus, weil $\delta\delta=1$ ist, leicht geschlossen wird, dass
\[
p^{\frac12(q-1)}-\gamma\delta
\]
durch $q$ theilbar sein müsse. Da nun auch $p^{\frac12(q-1)}-\beta$ durch $q$ nach einem bekannten Theorem theilbar ist, so wird nothwendig $\beta=\gamma\delta$ sein, woraus wiederum das Fundamental-Theorem von selbst folgt.

\SourceLeaf{174}%
Das Fundamental-Theorem, verbunden mit einigen bekannten Lehrsätzen, kann zwar zu einer ziemlich kurzen Auflösung der Aufgabe dienen, zu entscheiden, ob eine vorgegebene ganze positive Zahl von einer gegebenen Primzahl quadratischer Rest oder Nichtrest sei, wie in der Abhandlung ausführlich gezeigt ist. Allein bei weiterm Nachdenken über den dritten Beweis des Fundamental-Theorems kam der Verf. auf eine noch viel geschmeidigere Auflösung, welche die dritte Abtheilung der Abhandlung ausmacht, und wovon wir hier blos die Endregel hersetzen, indem wir die Entwickelung ihrer Gründe Kürze halber übergehen. Wenn entschieden werden soll, ob die ganze positive Zahl $b$, welche durch die Primzahl $a$ nicht theilbar ist, von dieser ein quadratischer Rest oder Nichtrest sei, so bilde man, ganz auf dieselbe Art, wie wenn der grösste gemeinschaftliche Divisor von $a$ und $b$ gesucht werden sollte, die Gleichungen
\[
\begin{aligned}
a&=\beta b+c\\
b&=\gamma c+d\\
c&=\delta d+e\\
d&=\varepsilon e+f\text{ u. s. w.}
\end{aligned}
\]
bis man in der Reihe der Zahlen $a,b,c,d,e,f$ u. s. w. auf .die Einheit kommt. Man. bezeichne die Zahlen $\frac12a,\frac12b,\frac12c,\frac12d$ u. s. w., mit Weglassung des ihnen anhängenden Bruches $\frac12$, in so fern einige der Zahlen $a,b,c,d$ u. s. w. ungerade sind, durch $a',b',c',d'$ u. s. w.; man nenne $\mu$ die Anzahl der in der Reihe $a',b',c',d'$ u. s. w. vorkommenden Folgen zweier ungeraden Zahlen unmittelbar nach einander, endlich nenne man $\nu$ die Anzahl derjenigen ungeraden Zahlen in der Reihe $\beta,\gamma,\delta,\varepsilon$ u. s. w., welchen in der Reihe $b',c',d',e'$ u. s. w. der Ordnung nach eine Zahl von der Form $4k+1$ oder $4k+2$ entspricht. Dieses vorausgesetzt, wird $b$ quadratischer Rest oder Nichtrest von $a$ sein, je nachdem $\mu+\nu$ gerade oder ungerade ist, den einzigen Fall ausgenommen, wo zugleich $b$ gerade und $a$ von der Form $8k+3$ oder $8k+5$ ist, in welchem von jener Regel das Gegentheil Statt findet, so dass ein gerades $\mu+\nu$ anzeigt, dass $b$ quadratischer Nichtrest von $a$ ist, ein ungerades $\mu+\nu$ hingegen, dass $b$ quadratischer Rest von $a$ ist.
\BookRule{.28\linewidth}
